\documentclass[utf8]{frontiers_suppmat}

\usepackage[hyphens]{url}
\usepackage{xr-hyper}
\usepackage[hypertexnames=false]{hyperref}
\externaldocument{CivicWorkOS}
\usepackage{array}
\usepackage{tabularx}
\usepackage{booktabs}
\usepackage{siunitx}
\sisetup{detect-all}
\usepackage{float}
\usepackage[dvipsnames]{xcolor}
\usepackage{tikz}
\usetikzlibrary{arrows.meta,positioning}
\usepackage{pgfplots}
\pgfplotsset{compat=1.18}
\usepackage{threeparttable}
\usepackage{multirow}
\usepackage{adjustbox}
\hypersetup{
	colorlinks=true,
	linkcolor=MidnightBlue,
	urlcolor=MidnightBlue,
	citecolor=MidnightBlue
}

\renewcommand{\theHtable}{S\arabic{table}}
\renewcommand{\theHfigure}{S\arabic{figure}}
\renewcommand{\thesection}{S\arabic{section}}

\def\firstAuthorLast{Syed et~al.}

\begin{document}

\title[CivicWorkOS Supplementary]{CivicWorkOS: Capability-Preserving Allocation of Municipal Work Among Humans, AI Agents and Robots}

\maketitle

\section{The Contestability Procedure in Full}
\label{app:contest}

This appendix is the operative specification of the procedure summarized in Section~\ref{sec:contest}. Where the body of the article states a window or a threshold, the value given here governs.

\subsection{Standing}
\label{app:contest-standing}

An appeal against an allocation decision may be lodged by:

\begin{enumerate}
	\item any resident materially affected by the service outcome the allocation produced;
	\item any worker whose allocation, access to budget-protected developmental work, or displacement is at issue, \emph{irrespective of contract status}, whether permanent, fixed-term, agency, contracted, or platform-mediated;
	\item an authorized representative of either, including a recognized trade union or a community organization acting on a resident's behalf; and
	\item the municipal ombudsman or equivalent office, acting on its own motion where a pattern of outcomes rather than a single decision is at issue.
\end{enumerate}

Clause 2 is deliberately decoupled from the competent-practitioner headcount $N_k(t)$ used to size the capability budget. A worker who is invisible to the budget arithmetic can still contest a decision that affects them. This does not repair the exclusion documented in Section~\ref{sec:whowins}, but it ensures the exclusion is at least reportable by the people it affects. Clause 4 is the systemic-challenge route: an individual appellant contests one decision, whereas the ombudsman may contest the configuration that produced a pattern of them.

\subsection{Windows}
\label{app:contest-windows}

Table~\ref{tab:app-windows} gives the operative appeal and resolution windows by oversight zone and rights impact.

\begin{table}[!htbp]
	\caption{Appeal windows by oversight zone and rights impact. All values are policy parameters set by the Weight Review Panel; those below are the defaults proposed in Section~\ref{sec:contest}, which reproduces them exactly.}
	\label{tab:app-windows}
	\small
	\begin{tabularx}{\textwidth}{@{}p{3.6cm}p{2.3cm}p{2.3cm}X@{}}
		\toprule
		\textbf{Decision class} & \textbf{Lodgement} & \textbf{Resolution} & \textbf{Effect of lodgement}\\
		\midrule
		Zone Z1 or Z2, routine & 20 working days & 20 working days & Outcome stands pending resolution\\
		Zone Z1 or Z2, rights-affecting ($auth_i$ or $priv_i$ above threshold) & 20 working days & 10 working days & Outcome suspended pending resolution\\
		Zone Z3 or Z4 & 20 working days & 20 working days & Outcome stands; the human decision-maker of record is named in the response\\
		Worker access or displacement appeal & 40 working days & 30 working days & Displacement suspended pending resolution\\
		\bottomrule
	\end{tabularx}
\end{table}

The longer lodgement window for worker appeals reflects the fact that a worker may not observe the pattern of allocations affecting their developmental access until a budget period has closed. The rights-affecting threshold on $auth_i$ and $priv_i$ defaults to 0.7 on the normalized scale of~\eqref{eq:task}.

\subsection{Evidentiary Record}
\label{app:contest-record}

Every allocation decision writes an append-only record containing:

\begin{enumerate}
	\item the task profile $T_i$ of~\eqref{eq:task}, including $d_i$ and the derived developmental content $\ell_i$;
	\item the mode set drafted, the roster set $\mathcal{A}_{i,m}$ constructed for each surviving mode, and the set surviving policy filtering, with the Policy Digital Twin status returned for each removed mode and each pruned roster, together with the identifier of every rule that fired;
	\item for every surviving staffed mode, the full term vector $(Q,S,P,Eq^{srv},Tr,Cost,En,Pr)$, the five debt components, the composed $CAD_{i,m,a}$, the base $SCV_{i,m,a}$, the Lagrangian-augmented score $\widetilde{SCV}_{i,m,a}$, and the credited developmental share $\phi_m\psi_a$;
	\item the admissibility outcome for each staffed mode, and where one was excluded, which bracket of the admissibility test excluded it and the value of the catch-up increment $\delta_k(t)$ in force;
	\item the dual prices $\lambda_k$, $\mu_s$, $\nu_{k,g}$, $\varrho_g$, $\upsilon_r$ in force and the timestamp of the rebalance that produced them, together with the integrality gap that rebalance reported;
	\item the debt-weight flip threshold $w_9^\star$ of Proposition~\ref{prop:price} for this task, so that an appellant can see whether the objective or a constraint determined the outcome;
	\item the version identifier of the parameter configuration, covering weights, budgets, access targets, the group coarsening and minimum cell size, transition ceilings, zone thresholds, and contingency parameters, and the panel decision that authorized it; and
	\item the identity and role of the accountable human where the zone designates one.
\end{enumerate}

The record is disclosable in full to an appellant with standing, subject only to redaction of third-party personal data and to the small-cell suppression rule of Section~\ref{sec:problem}, which prevents the record from identifying individual members of a group cell below $n^{min}$. Items 5 through 7 are what make an appeal about a \textit{policy} rather than about a single decision possible: an appellant who believes the weight configuration itself is unjustified can identify and cite the panel decision that set it, and item 6 tells them whether disputing the weights would have changed anything.

\subsection{Remedy and Systemic Trigger}
\label{app:contest-remedy}

An upheld appeal entitles the appellant to re-decision of the task class under Zone Z3, with a designated human decision-maker who is named in the response and who may depart from the mechanism's recommendation on the record.

Beyond individual remedy, if the upheld-appeal rate for a task class exceeds 0.05 within the rolling window $W$, the class is automatically demoted one oversight zone by predicate 2 of Table~\ref{tab:zonetrans} (Z1$\to$Z2, Z2$\to$Z3, Z3$\to$Z4) and referred to the Weight Review Panel at its next sitting. The demotion is recorded as a versioned policy update in the same way as any other panel decision, so the appeal history that triggered it remains traceable. Restoration to the previous zone requires an affirmative panel decision and cannot occur automatically.

\section{The Specified Evaluation Protocol}
\label{app:protocol}

This appendix records in full the evaluation protocol that Section~\ref{sec:protocol} of the main article summarizes. \textbf{The protocol has not been executed.} It is set out at this level of detail so that the evaluation the framework requires is reproducible by others and auditable against a plan fixed in advance rather than against one reconstructed afterwards. No quantity anywhere in this appendix is an observed outcome. The one column of measured-looking numbers, the ``worked case'' column of Table~\ref{tab:app-predictions}, is the exact arithmetic of Section~\ref{sec:worked} of the main article, which is what the refutation thresholds were calibrated against.

\subsection{Testbed and Sectors}
\label{app:protocol-testbed}

The protocol specifies an urban digital twin and multi-agent discrete-event simulation testbed across six municipal sectors, with task arrivals from synthetic but calibrated streams over a ten-year horizon so that slow-moving effects such as retirement waves and pipeline decay would become observable. Table~\ref{tab:app-sectors} gives the sector scale and calibration status. The city to be modeled is a metropolitan authority of roughly 1.5 million residents, and the annual task volume of 803{,}100 sets the scale at which the tractability of the periodic program~\eqref{eq:program} has to be assessed.

\begin{table}[!htbp]
	\caption{Evaluated sectors, annual task volumes, and calibration status. ``Calibrated'' means arrival rates and volume distributions are fitted to the named open municipal series. ``Parametric'' means no comparable open series was identified and the sector's distributions are specified as a swept range rather than fitted, with results reported separately.}
	\label{tab:app-sectors}
	\small
	\begin{tabularx}{\textwidth}{@{}lrXl@{}}
		\toprule
		\textbf{Sector} & \textbf{Tasks/yr} & \textbf{Calibration source} & \textbf{Status}\\
		\midrule
		Infrastructure inspection & 2{,}100 & Asset-condition and inspection-interval records from the portals below and published national bridge inventories & Calibrated\\
		Waste collection and sanitation & 186{,}000 & NYC Open Data, Department of Sanitation monthly tonnage and collection-route geography & Calibrated\\
		Citizen-service administration & 412{,}000 & NYC Open Data 311 service requests; Open Data BCN incident and service-request series & Calibrated\\
		Public-transport operations & 94{,}000 & London Datastore, Transport for London performance and ridership releases & Calibrated\\
		Emergency logistics & 31{,}000 & No comparable open municipal series identified & Parametric\\
		Municipal facility management & 78{,}000 & Helsinki Region Infoshare municipal property and building registers & Calibrated\\
		Workforce composition & --- & Establishment and headcount tables of the same municipalities & Calibrated\\
		\bottomrule
	\end{tabularx}
\end{table}

Two calibration caveats belong with the protocol rather than buried in it. Emergency logistics has no calibration source. Earlier presentations of this protocol attributed its fleet-utilization distributions to two survey papers by Meseguer-Valenzuela and co-authors and by Hady and co-authors, which is incorrect: both are reviews of allocation methods and neither contains utilization distributions. The sector is therefore specified parametrically over a swept range of arrival rate and fleet utilization, its outcomes are to be reported separately rather than folded silently into the aggregate, and no calibration claim is made for it. Second, all five named portals belong to high-income cities with mature open-data programs, which under-represents exactly the municipalities the automation-exposure literature identifies as most vulnerable. This is not a conservative bias, as an earlier version of this article claimed. It runs the wrong way: the framework's central weakness is that it cannot see non-payroll workers, and the calibration cities are the ones with fewest of them, so the calibration would systematically understate the size of the population the framework fails. Appendix~\ref{app:calib} (Table~\ref{tab:app-calib}) records each series identifier and the transformation to be applied.

\subsection{Baseline Allocation Strategies}
\label{app:protocol-baselines}

Six strategies are to be compared. Four are configurations of~\eqref{eq:scv} and its constraints, and two are external methods from the human--robot collaboration literature, to be reimplemented against the same task representation.

\textit{Automation-First} maximizes quality and productivity, preferring automation wherever policy permits, with $w_9 = 0$ and constraints~\eqref{eq:prog-hcpb} through~\eqref{eq:prog-res} relaxed. \textit{Cost/Performance Optimization} maximizes $(Q,P)$ net of $Cost$ and $En$ only. \textit{Human-First} restricts $m$ to the admissible mode with the largest human share $\phi_m$ whenever a qualified human is available, and selects rosters on availability alone. \textit{Capability Matching} implements the capability-fit criterion of Ranz and co-authors, scoring each mode by the match between the task requirement vector and the mode's capability profile, with no debt, budget, access, or reserve terms. \textit{Ergonomics-Aware Role Allocation} implements the AND/OR-graph role allocation with ergonomic risk estimation of Merlo and co-authors, redirecting high-$phy$ and high-$risk$ actions away from human performers. \textit{CivicWorkOS} applies the full objective under all constraints of~\eqref{eq:program}.

Both external methods were designed for a workcell rather than a city, and reimplementing them requires two adaptations we state rather than bury: their capability and ergonomic profiles must be mapped onto the nine dimensions of~\eqref{eq:task}, and neither method selects a roster, so both would be given the availability-ordered roster that Human-First uses. Neither adaptation advantages CivicWorkOS, and both are to be documented in the replication package so that the mapping can be challenged.

The limit of this design should be stated where the design is, and not only where its results would be. An evaluation in which four of six comparators are configurations of the proposed objective can show that the mechanism does what it was designed to do; it cannot show that it outperforms an independently designed alternative. The two external baselines are specified for that reason, and the comparison is incomplete until they are implemented and run. Neither is implemented at the time of writing.

\subsection{Metrics}
\label{app:protocol-metrics}

Fourteen metrics are to be tracked, of which twelve are comparable across strategies and enter the hypothesis family, and two are to be reported descriptively because they have no better-or-worse orientation for a strategy that lacks the mechanism producing them.

Six comparable metrics are operational: \textit{Productivity} and \textit{Operating Cost} indices normalized to Automation-First at 100, \textit{Service Quality} on a 0--100 scale, \textit{Safety Incidents} per 10{,}000 tasks, an \textit{Energy Use} index, and \textit{Mean Recovery Time}, hours to restore $\kappa_s$ of critical-service capacity after an acute disruption. Two concern capability: the \textit{Capability-Formation Index}, $\Lambda_k/B_k$ summed over domains, and \textit{Accumulated CAD}, the running sum of~\eqref{eq:cad} normalized to 100 for Automation-First at year ten. One concerns the mechanism's own failure mode: the \textit{Zone Z4 fallback rate}, the share of tasks for which $\mathcal{A}(i,t)=\emptyset$ and the task routes to human-reserved execution. A framework with more constraints empties more candidate sets, and reporting this is how a reader would see the price of the constraint system in operational terms rather than only in objective units.

Three comparable metrics are distributional: \textit{protected-practice share by group}, $\Pi_{k,g}$ from~\eqref{eq:accessshare}, to be reported per group and summarized by the share reaching the least-represented group; \textit{displacement by group}, $\Delta_g$, summarized by the most-affected group; and \textit{service-equity dispersion}, the coefficient of variation of realized service quality across resident districts, which is what $Eq^{srv}$ aggregates and a single equity scalar hides. All three are to be reported as group-level vectors rather than only as the summary statistic, because a framework that answers who wins and who loses with one number has not answered it.

Two metrics are descriptive. The \textit{realized dual} $\bar\lambda$ is to be reported because it is the framework's own statement of what preservation costs, and it is identically zero for every strategy that does not impose the constraint, so comparing it across strategies is meaningless. \textit{Contestability throughput}, appeals lodged, upheld, and resolved within the windows of Section~\ref{sec:contest}, exists only for CivicWorkOS because the other strategies have no appeals procedure. Excluding these two from the hypothesis family repairs a defect in earlier versions of this protocol, where a hypothesis asserting that every strategy is worst on at least one of thirteen metrics was undefined for metrics that admit no ordering.

Reporting contestability throughput at all requires a generative model of appeals that this protocol does not yet contain. Arrival of appeals, determination of standing, the upheld rate, and the resolution-time distribution would each have to be specified and justified before any throughput figure could be reported, and the systemic trigger of Section~\ref{app:contest-remedy} depends on all four. Until that model exists, the metric is named in the protocol and no value for it is reportable.

\subsection{Implementation Configuration and Default Parameters}
\label{app:protocol-config}

The implementation pairs a Python discrete-event simulation of task arrival, roster construction, agent pricing, and execution with a mixed-integer solver invoked at the rebalancing cadence to solve~\eqref{eq:program}, since the online rule does not by itself guarantee the coupled constraints. Reference hardware is a single 32-core, 128\,GB workstation-class server. Earlier versions of this configuration specified hyperparameters for an optional learned surrogate that was never used or evaluated; it has been removed, and no learned component appears anywhere in the framework.

Because~\eqref{eq:program} is a binary multi-dimensional assignment problem with side constraints, the rebalance window is the parameter that governs tractability, and the protocol requires the solve time, binary-variable count and optimality gap to be reported at each cadence rather than asserted. The weekly default, the monthly, quarterly and annual alternatives, and a sector-decomposed annual variant that trades the cross-sector coupling of~\eqref{eq:prog-cap} for tractability are all to be profiled. No such profile has been measured, and the main article makes no scalability claim on this basis.

Table~\ref{tab:app-params} assigns a default to every remaining parameter in the framework. The objective and debt weights themselves are given in Section~\ref{sec:weights} of the main article.

\begin{table}[!htbp]
	\caption{Default values for every framework parameter. Domain-specific values are shown for structural inspection where they differ. Panel-set parameters are those the Weight Review Panel of Section~\ref{sec:feedback} owns; the rest follow from continuity planning, certification, or numerical requirements.}
	\label{tab:app-params}
	\small
	\begin{tabularx}{\textwidth}{@{}llXl@{}}
		\toprule
		\textbf{Parameter} & \textbf{Symbol} & \textbf{Default} & \textbf{Owner}\\
		\midrule
		Budget period & $\Delta T$ & 1 year & Panel\\
		Replacement factor & $\eta_k$ & 1.2 & Panel\\
		Budget floor & $B^{min}_k$ & 2{,}000 h (inspection); $0.05\,\Phi_k$ elsewhere & Panel\\
		Raw competence hours & $h^{raw}_k$ & 1{,}800 h (inspection); per certification elsewhere & Certification\\
		Access tolerance & $\varepsilon_k$ & 0.05 & Panel\\
		Minimum group cell size & $n^{min}$ & 30 practitioners per cell per domain & Panel\\
		Displacement ceiling & $\tau_g$ & 0.10 of baseline hours per budget period & Panel\\
		Reskilling coverage floor & $\chi^{min}$ & 0.80 & Panel\\
		Access slack penalty & $M$ & $10^{3}$ objective units per unit share & Numerical\\
		Contingency order & $n_s$ & 2 for emergency logistics and transport; 1 otherwise & Continuity plan\\
		Peak-demand share & $\kappa_s$ & 1.00 emergency logistics; 0.80 transport and sanitation; 0.60 otherwise & Continuity plan\\
		Safety floor & $S^{min}_i$ & $risk_i\cdot crit_i$ & Safety Agent\\
		Zone Z3 human-share floor & $\phi^{Z3}$ & 0.50 & Panel\\
		Zone evaluation window & $W$ & 4 budget quarters & Panel\\
		Promotion appeal-rate ceiling & --- & 0.02 upheld over $W$ & Panel\\
		Demotion appeal-rate trigger & --- & 0.05 upheld over $W$ & Panel\\
		Drift trigger & --- & Population Stability Index (PSI) $>0.25$ & Standard practice\\
		Rebalance cadence & --- & weekly & Operational\\
		Trainee domain hours & $\Theta_k$ & 1{,}500 clock hours per year & Establishment tables\\
		\bottomrule
	\end{tabularx}
\end{table}

\subsection{Statistical Analysis Plan}
\label{app:protocol-stats}

Each strategy is to be run for at least 30 independent seeds per sector, each seed governing one complete ten-year run. Seeds are to be drawn once and reused across strategies, so that comparisons are paired on the arrival stream, roster availability, negotiation timing, and failure-injection schedule. Outcomes are to be reported as means with 95\% confidence intervals, or bootstrap percentile intervals over 10{,}000 resamples where the sampling distribution fails a Shapiro--Wilk test at $\alpha=0.05$ or the metric is bounded and its mean lies within one standard deviation of a bound.

Pairwise comparisons between CivicWorkOS and each of the five baselines use the \textit{Wilcoxon signed-rank test}. Earlier versions of this plan specified the Mann--Whitney $U$ test, which is a test for independent samples and is inconsistent with a paired-seed design; the signed-rank test is the correct paired non-parametric analogue, and it is chosen over the paired $t$-test because several metrics, including safety incidents per 10{,}000 tasks and recovery time, are bounded below and expected to be right-skewed. The family for multiplicity correction is the full grid of 12 comparable metrics $\times$ 5 strategy comparisons $=$ 60 tests, to be controlled by the Holm--Bonferroni step-down procedure at a family-wise $\alpha=0.05$.

The smallest per-comparison $\alpha$ in that family is $0.05/60 \approx 0.000833$, giving a two-tailed critical $z$ of 3.34. With $z_\beta = 0.84$ for 80\% power, the minimum detectable standardized effect over 30 paired observations is $d = (3.34+0.84)/\sqrt{30} = 0.76$, and the signed-rank test retains roughly 95\% of that efficiency for the distributions we expect. The normal approximation is mildly optimistic and the $t$ reference distribution pushes the requirement slightly higher, so 0.76 should be read as a floor rather than an exact figure. Where the refutation threshold registered for a metric in Table~\ref{tab:app-predictions} implies a smaller standardized effect than 0.76, the seed count for that metric is to be raised until the threshold is detectable at 80\% power, and both the raised seed count and the achieved power are to be reported alongside the outcome. This rule is fixed in advance so that seed counts cannot be chosen after inspecting outcomes.

Distributional metrics are to be reported per group without aggregation across groups. No significance claim is to be made for a group represented by fewer than $n^{min}=30$ simulated practitioners, consistent with the group-definition rule of Section~\ref{sec:problem}; such groups are to be reported descriptively with their sample size stated. Any deviation from this plan, such as a changed test, a changed family, or a changed seed count not covered by the rule above, is to be documented in the replication audit log with the rationale and timestamp of the determination.

\subsection{Stress-Test Protocol}
\label{app:protocol-stress}

Seven scenarios address RQ6 and RQ7. Acute scenarios inject a single disruption in separate runs to avoid confounding, and chronic scenarios run for the full horizon. Injection magnitudes are in Appendix~\ref{app:stress} (Table~\ref{tab:app-stress}). Scenarios (i) through (iv) are an AI-service outage, a cyberattack on fleet command and control, a robot-fleet mechanical failure, and an emergency demand surge, each to be evaluated by hours to restore $\kappa_s D^{peak}_s$. Scenarios (v) through (vii) are a retirement wave, rapid improvement in AI capability, and a new municipal regulation entering the Policy Digital Twin without redeploying the market, to be evaluated by the trajectories of $\Lambda_k$, $\lambda_k$, and $\Pi_{k,g}$, by zone-transition counts under the predicates of Table~\ref{tab:zonetrans}, and by constraint-update latency.

Scenario (v) is the framework's self-test and is to be run in two arms, because Proposition~\ref{prop:feasibility} shows that the single-arm design used in earlier versions of this protocol was arithmetically impossible. Arm (v-a) raises $r_k$ to 0.14, inside the critical rate of 0.1418 for structural inspection, and tests whether the capability feedback is stabilizing: raising $r_k$ raises $B_k$ through~\eqref{eq:hcpb-estimator}, hence $\lambda_k$, pushing allocation toward developmental rosters at exactly the moment fewer competent leads are available to supervise them. Arm (v-b) raises $r_k$ to 0.25, which puts $B_k = 10{,}368$ hours against a domain ceiling of $\Phi_k\bar\zeta_k = 5{,}880$ and is therefore infeasible by construction. Arm (v-b) does not test stability, which is undefined in that regime. It tests the escalation path of Section~\ref{sec:feasibility}: whether the infeasibility flag raises, whether the deficit is recorded correctly, and whether the four restoration options are reported with the right magnitudes. Running an infeasible parameter set as though it were a stability test, as an earlier version of this protocol did, would have fired a refutation criterion by arithmetic before any dynamics ran. Section~\ref{sec:shocks} of the main article gives the exact arithmetic of both arms, which is the target the implementation would have to reproduce.

\subsection{Hypotheses and Refutation Criteria}
\label{app:protocol-hypotheses}

Table~\ref{tab:app-predictions} specifies ten hypotheses with explicit refutation criteria. Three properties distinguish this set from the one in earlier versions of this article. Every threshold has been checked against the worked example of Section~\ref{sec:worked} and, where the single-domain case falls outside it, either the threshold has been revised with the reason stated or the divergence is recorded in the hypothesis itself. P8 and P10 have been restated so that they are falsifiable rather than guaranteed by construction. And P1's refutation criterion no longer uses overlapping confidence intervals, which is not a valid test, since non-overlapping intervals imply significance but the converse does not hold.

These hypotheses are \textit{pre-specified}, not pre-registered. Pre-registration means a timestamped third-party registry entry made before data collection, and stating hypotheses in a manuscript is not that. The set below will be lodged on the Open Science Framework before the evaluation is executed, and the registration identifier will be cited in whichever article reports the outcome. Until then the correct word is pre-specified, and we use it. No hypothesis in Table~\ref{tab:app-predictions} has been tested, and the table records no verdicts.

\begin{table}[!htbp]
	\caption{Pre-specified evaluation hypotheses and refutation criteria. The final column reports the value the single-domain worked example of Section~\ref{sec:worked} produces, which is what the thresholds were calibrated against. A dash means the quantity is not defined for a single domain. No hypothesis in this table has been tested; the protocol specifying the test has not been executed.}
	\label{tab:app-predictions}
	\footnotesize
	\begin{tabularx}{\textwidth}{@{}p{0.5cm}XXp{1.9cm}@{}}
		\toprule
		\textbf{ID} & \textbf{Prediction} & \textbf{Refuted if} & \textbf{Worked case}\\
		\midrule
		P1 & Accumulated CAD at year 10 under CivicWorkOS is below half that under Automation-First. & The ratio is $\geq 0.5$, or the Holm-corrected signed-rank test does not reject at family-wise $\alpha=0.05$. & 0.79 (single task class, no accumulation)\\
		P2 & CivicWorkOS retains at least 80\% of Automation-First productivity. & Productivity index $< 80$ against Automation-First at 100. & 88.5\\
		P3 & The Capability-Formation Index under CivicWorkOS is at least 0.95, and exceeds that of Capability Matching by at least 40 index points on a 0--100 scale. & CFI $< 0.95$, or the gap is below 40 points. & 1.00 against 0.00\\
		P4 & Mean recovery time under CivicWorkOS is at most half that under Automation-First across scenarios (i)--(iv). & The ratio exceeds 0.5 in two or more of the four scenarios. & ---\\
		P5 & Disabling the HCPB raises accumulated CAD more than disabling any other single mechanism. & Any other single-mechanism ablation raises CAD by more, by a margin exceeding 5 index points. & ---\\
		P6 & Disabling the resilience reserve at least doubles mean recovery time while changing accumulated CAD by less than 20\%, establishing that resilience and debt are not interchangeable. & Recovery time less than doubles, or CAD moves by 20\% or more. & ---\\
		P7 & In arm (v-a), $\lambda_k$ rises and allocation shifts toward developmental rosters without oscillation, recovering $\Lambda_k \geq B_k$ within three budget periods. In arm (v-b), the infeasibility flag raises within one budget period and the recorded deficit matches $B_k - \Phi_k\bar\zeta_k$ to within 1\%. & $\lambda_k$ or the developmental-roster share oscillates with increasing amplitude in (v-a), or (v-a) fails to recover within three periods, or (v-b) fails to flag or misreports the deficit. & (v-b) deficit $=4{,}488$ h by Proposition~\ref{prop:feasibility}\\
		P8 & Disabling constraint~\eqref{eq:access} leaves $\Pi_{k,g}$ for the least-represented group within 0.02 of the incumbent share in every sector; enabling it raises $\Pi_{k,g}$ by at least 0.12 while raising accumulated CAD by no more than 5 index points. & Disabling it moves $\Pi_{k,g}$ more than 0.02 from the incumbent share, or enabling it raises $\Pi_{k,g}$ by less than 0.12, or the CAD cost exceeds 5 points. & $+0.145$ at zero CAD cost\\
		P9 & CivicWorkOS operating cost does not exceed 130\% of Automation-First. & Cost index $> 130$. & 124.5\\
		P10 & No strategy is weakly best on all twelve comparable metrics. & Any strategy is weakly best on all twelve. & CivicWorkOS worst on cost\\
		\bottomrule
	\end{tabularx}
\end{table}

Two thresholds were revised against the worked example rather than retained. P9's ceiling was raised from 125\% to 130\% because the worked optimum costs 124.5\% of Automation-First in a single domain, and a threshold that the framework's own demonstration case clears by half a percentage point is not a threshold, it is a coincidence. P3's ratio criterion was replaced by an absolute floor plus a gap, because the ratio is undefined whenever the comparator delivers zero developmental practice, which is exactly what Capability Matching does in the worked domain. Both revisions are recorded here so that a reader can see the thresholds were set with the arithmetic in hand rather than before it. Setting a threshold against arithmetic one already holds is a weaker form of pre-specification than setting it against data one does not, and we claim no more than the weaker form.


\section{Stress-Test Injection Schedules}
\label{app:stress}

This appendix gives the injection magnitudes for the seven scenarios of Appendix~\ref{app:protocol-stress} (Table~\ref{tab:app-stress}). Acute scenarios (i) through (iv) are injected independently in separate runs to avoid confounding, and chronic scenarios (v) through (vii) run for the full horizon.

\begin{table}[!htbp]
	\caption{Stress scenarios, injection points, and evaluation criteria. Scenario (v) is split into two arms because Proposition~\ref{prop:feasibility} shows that a single-arm design at $r_k=0.25$ is arithmetically infeasible for the structural-inspection domain and therefore cannot test stability.}
	\label{tab:app-stress}
	\small
	\begin{tabularx}{\textwidth}{@{}p{0.6cm}p{3.5cm}XX@{}}
		\toprule
		& \textbf{Scenario} & \textbf{Injection} & \textbf{Evaluated by}\\
		\midrule
		i & AI-service outage & 100\% of AI-agent capacity removed for 72 simulated hours at year 5 & Hours to restore $\kappa_s D^{peak}_s$\\
		ii & Cyberattack on fleet command and control & 100\% of robot-fleet capacity removed for 168 hours; policy twin raises affected classes to Z3 & Hours to restore $\kappa_s D^{peak}_s$; zone transitions\\
		iii & Robot-fleet mechanical failure & 40\% of one fleet class removed for 30 days & Hours to restore $\kappa_s D^{peak}_s$\\
		iv & Emergency demand surge & Arrival rate $\times 3$ for 14 days in the emergency-logistics and facility sectors & Hours to restore $\kappa_s D^{peak}_s$\\
		v-a & Retirement wave, feasible arm & $r_k$ raised from 0.12 to 0.14 for two consecutive budget periods, inside the critical rate $r^{crit}_k = 0.1418$ & Trajectories of $B_k$, $\lambda_k$, $\Lambda_k$, developmental-roster share; oscillation check per P7\\
		v-b & Retirement wave, infeasible arm & $r_k$ raised to 0.25, exceeding $r^{crit}_k$ by 76\% & Infeasibility flag latency; recorded deficit against $B_k-\Phi_k\bar\zeta_k$; correctness of the four reported restoration options\\
		vi & Rapid AI capability improvement & AI-agent $Q$ and $Tr$ raised 20\% over 3 years & Z2$\to$Z1 transition counts under the promotion predicates of Table~\ref{tab:zonetrans}\\
		vii & New municipal regulation & A new \textit{prohibit} rule added to the Policy Digital Twin at year 6 without redeploying the market & Constraint-update latency; allocation shift; consistency-check outcome\\
		\bottomrule
	\end{tabularx}
\end{table}

Scenario (v) is the framework's self-test. Raising $r_k$ raises $B_k$ through~\eqref{eq:hcpb-estimator}, which raises the dual price $\lambda_k$ and pushes allocation toward developmental rosters at precisely the moment fewer competent leads are available to supervise them. Prediction P7 of Table~\ref{tab:app-predictions} defines the refutation criterion for each arm. Arm (v-a) is refuted by oscillation with increasing amplitude or by failure to recover $\Lambda_k \geq B_k$ within three budget periods. Arm (v-b) is refuted if the infeasibility flag does not raise within one budget period or if the recorded deficit departs from $B_k - \Phi_k\bar\zeta_k$ by more than 1\%. Stability is undefined in the (v-b) regime and is not tested there.

\section{Calibration Sources and Transformations}
\label{app:calib}

This appendix expands the calibration column of Table~\ref{tab:app-sectors} with the series used and the transformation applied to each (Table~\ref{tab:app-calib}). No archived snapshots accompany this article. Each series is named with its issuing body and the transformation applied to it, so that a reader can retrieve the current release directly from the source; we do not report retrieval timestamps, because a timestamp without the corresponding snapshot is not verifiable. These figures size the sectors of a protocol that has not been executed, and no result reported in this article depends on any of them.

\begin{table}[!htbp]
	\caption{Calibration sources, series, and transformations across the evaluated sectors. Emergency logistics has no calibration source and is specified parametrically; no calibration claim is made for it.}
	\label{tab:app-calib}
	\footnotesize
	\setlength{\tabcolsep}{3pt}
	\begin{tabularx}{\textwidth}{@{}p{2.3cm}p{2.4cm}XX@{}}
		\toprule
		\textbf{Sector} & \textbf{Portal} & \textbf{Series} & \textbf{Transformation applied}\\
		\midrule
		Waste collection and sanitation & NYC Open Data & Department of Sanitation monthly tonnage; collection-route and district geography & Monthly tonnage converted to route-day task counts by district, then to a non-homogeneous Poisson arrival rate with weekly and seasonal components\\
		Public-transport operations & London Datastore & Transport for London performance and ridership releases & Performance-incident counts converted to maintenance and intervention task arrivals; ridership used to set $D^{peak}_s$\\
		Citizen-service administration & NYC Open Data; Open Data BCN & 311 service requests; Barcelona incident and service-request series & Request records grouped into task classes by complaint type, with service-level targets giving $urg_i$ and closure times giving $d_i$\\
		Municipal facility management & Helsinki Region Infoshare & Municipal property and building registers & Building stock and asset class converted to periodic maintenance and inspection task streams on statutory intervals\\
		Infrastructure inspection & Open portals of the above municipalities; published national bridge inventories & Asset condition and inspection-interval records & Structure counts and statutory inspection intervals converted to annual inspection events; team sizes from establishment tables\\
		Emergency logistics & No comparable open series identified & --- & Arrival rate and fleet utilization specified parametrically and swept; results reported separately and not aggregated as calibrated\\
		Workforce composition & Establishment and headcount tables of the same municipalities & Practitioner counts by domain; age and tenure distributions where published & Headcount by domain gives $N_k$; age and tenure distributions give $r_k$ by cohort projection\\
		\bottomrule
	\end{tabularx}
\end{table}

The direction of the sampling bias in this design is stated in Appendix~\ref{app:protocol-testbed}, and it runs against the framework rather than for it. One further limit belongs here. Open portals do not publish practitioner-level developmental logs, so $learn_i$ and the derived $\ell_i$, the mode shares $\phi_m$, the supervision ratios that bound $\psi_a$, and the certification hours $h^{raw}_k$ cannot be calibrated from them. Appendix~\ref{app:elicit} specifies the elicitation protocol for those parameters. That protocol has \emph{not} been executed. The values used throughout this article are the authors' own structured estimates, and Section~\ref{sec:sensitivity} maps the region over which each conclusion survives them.

\section{Elicitation Protocol for the Judgment Parameters}
\label{app:elicit}

Four parameters in CivicWorkOS are structured professional judgments rather than measurements: the task learning value $learn_i$, the human developmental share $\phi_m$ of each hybrid mode, the supervision ratio that bounds the roster eligibility $\psi_a$, and the certification requirement $h^{raw}_k$. The first three carry the framework's headline result and the fourth determines whether the capability budget is feasible at all, so this appendix specifies the instrument that would replace the estimates with data. Nothing in this appendix has been carried out.

\subsection{Participants and Qualification}
\label{app:elicit-participants}

Per domain, a minimum of 12 raters drawn from three groups in equal proportion: practitioners currently competent in the domain with at least five years of practice, practitioners currently developing in the domain, and the supervisors responsible for signing off developmental progression. Including developing practitioners is not a courtesy. They are the only group with direct experience of what a machine-assisted task teaches, and excluding them would reproduce inside the instrument the same incumbent bias that Section~\ref{sec:access} identifies in the budget.

\subsection{Instrument}
\label{app:elicit-instrument}

Each parameter is rated on an anchored eleven-point scale with behaviourally defined endpoints and a defined midpoint.

For $learn_i$, the anchors are: 0, a task from which a developing practitioner acquires no judgment transferable to any other task in the domain; 0.5, a task from which a developing practitioner acquires judgment transferable within the same task class only; 1.0, a task from which a developing practitioner acquires judgment transferable across the domain and required for independent competence.

For $\phi_m$, raters are shown a task specification and a mode description and asked what share of the judgment the task requires is still exercised by the human, with anchors at 0 (the human observes an outcome and cannot alter it), 0.5 (the human sets the parameters of a machine determination and reviews its output against independent evidence), and 1.0 (the human forms the determination and the machine supplies data only).

For the supervision ratio, raters are asked the maximum number of developing practitioners one competent lead can supervise on the given task class without the lead's own attention to the task being compromised, conditioned on the task's $risk_i$ and $crit_i$.

For $h^{raw}_k$, raters supply the supervised clock hours the certification route requires, which is a documentary question, and separately their estimate of the hours actually required in practice, which is a judgment. Divergence between the two is itself reportable.

\subsection{Procedure and Reliability}
\label{app:elicit-procedure}

Ratings are collected in two rounds on a Delphi pattern. Round one is independent and unseen. Round two shows each rater the anonymized distribution of round-one ratings and the free-text rationales attached to the extreme ratings, and invites revision. Agreement is assessed on round-two ratings by the two-way random-effects intraclass correlation coefficient ICC(2,k) with a pre-set acceptance threshold of 0.75. Where agreement falls below the threshold for a parameter, that parameter is reported as a range rather than a point value, and the sensitivity analysis of Section~\ref{sec:sensitivity} is re-run across the observed range rather than the illustrative one.

\subsection{Ethics and Data Protection}
\label{app:elicit-ethics}

The protocol involves human participants giving professional judgments about their own work and their colleagues' development, so it requires institutional ethics approval and written informed consent before any data collection. Ratings are held pseudonymously, group-level results only are reported, and no rating is disclosable to a participant's employer in a form that identifies the rater. Because the ratings concern developmental progression, which affects the participants' own access to protected work, the consent form must state that participation is voluntary and that a refusal has no effect on the participant's allocation under any deployed system.

\section{Extended Related Work}
\label{app:rw}

This appendix gives the full survey that Section~\ref{sec:relatedwork} condenses. The main text states the positioning; what follows states it work by work, so that a reader checking whether a particular literature has been engaged, and on what terms, can do so without reconstructing it from the compressed account.

\subsection{Human--Robot and Human--AI Task Allocation}
\label{app:rw-hrta}

Ranz et al.~\cite{Ranz2017} proposed capability-based task allocation by matching task requirements to human and robot capabilities, establishing that allocation should depend on task--agent fit rather than on automation for its own sake, and Dorri et al.~\cite{Dorri2018} situate this within multi-agent systems theory, surveying the coordination and negotiation mechanisms that later allocation-market designs, including ours, build upon. Subsequent work expanded the criteria. El Makrini et al.~\cite{Makrini2019} considered ergonomics in collaborative assembly. Merlo et al.~\cite{Merlo2023} proposed dynamic role allocation using AND/OR graphs and ergonomic risk estimation, redirecting high-risk actions away from humans. Noormohammadi-Asl et al.~\cite{Noormohammadi2025} incorporated human leading and following preferences into adaptive robot planning, and Ali et al.~\cite{Ali2022} model artificial trust explicitly, re-optimizing allocation as measured trust evolves. Meseguer-Valenzuela and Blanes Noguera~\cite{MeseguerValenzuela2025} review allocation for mobile robot fleets, close to the sanitation, inspection, and logistics fleets a city must coordinate, and Hady et al.~\cite{Hady2025} survey multi-agent reinforcement learning for resource allocation. At city scale, Zhao et al.~\cite{Zhao2023} allocate spatio-temporal sensing tasks across a mobile crowd using a knowledge graph, which is the closest published treatment of municipal task allocation as a graph-structured assignment problem.

These contributions target an operational team and optimize performance, productivity, ergonomics, preferences, or trust. They do not ask whether repeated local decisions erode the human experience needed to sustain future expertise, emergency substitution capacity, or resilience, nor, when that erosion occurs, whose expertise is eroded first.

\subsection{Workforce Scheduling with Skill and Training Constraints}
\label{app:rw-sched}

The operations-research literature has treated constrained training-hour allocation for fifty years, and the capability budget of Section~\ref{sec:hcpb} is a member of that family rather than a departure from it. Burke et al.~\cite{Burke2004} survey nurse rostering, where skill mix, supervision ratios, and qualification maintenance appear as hard constraints on a staffing decision, and Van den Bergh et al.~\cite{VandenBergh2013} review personnel scheduling more broadly, including the cross-training and skill-acquisition variants that decide how many staff reach a given competence by a given date.

Two differences matter for a municipal allocation mechanism. Rostering assumes the work exists and asks who covers it. Our decision is prior: whether the work is performed by a human at all, and therefore whether the developmental content exists to be allocated. And rostering optimizes coverage subject to skill constraints, whereas the objective here prices the future value of the practice itself. What we take from this literature is its discipline about supervision ratios and qualification accounting, which is exactly what Section~\ref{sec:hcpb} repairs relative to earlier presentations of this framework.

\subsection{Automation, Skill Erosion, and Capability Formation}
\label{app:rw-skill}

The concern at the center of this paper is not new, and we want its lineage on the record before claiming anything. Bainbridge~\cite{Bainbridge1983} identified the ironies of automation: automating the easy parts of a task degrades operator skill through disuse and leaves the human unable to intervene competently at exactly the moment intervention is required. Endsley and Kiris~\cite{EndsleyKiris1995} formalized the resulting out-of-the-loop performance problem, and Parasuraman and Riley~\cite{ParasuramanRiley1997} organized the failure modes into use, misuse, disuse, and abuse. Sheridan and Verplank~\cite{SheridanVerplank1978} supplied the levels-of-automation scale, generalized by Parasuraman et al.~\cite{Parasuraman2000} across information acquisition, analysis, decision, and action. Shneiderman~\cite{Shneiderman2020} recasts the trade-off as two independent axes of automation and human control, and Zuboff~\cite{Zuboff1988} draws the organizational distinction between automating and informating work, which is the distinction between a task that displaces a practitioner and one that develops them.

The quantitative basis for calibrating skill loss also exists. Arthur et al.~\cite{Arthur1998} meta-analyzed skill decay over non-use intervals across task types, and Tatel and Ackerman~\cite{TatelAckerman2025} update that estimate for procedural skills specifically, reporting retention as a function of retention interval, task closed-loop character, and rehearsal. These are the instruments a city would use to calibrate $D^{skill}$ against observed practice gaps rather than against expert judgment alone.

Three consequences follow. The mechanism of skill erosion is established, so CivicWorkOS must not claim to have discovered it. That literature is written at the scale of the individual operator, whereas the municipal question, whether a city's whole expertise pipeline is drawn down across thousands of independent decisions in dozens of services, is what we address. And it warrants treating fallback capacity as distinct from present competence, which is why $D^{fall}$ is separated from $D^{skill}$ in Section~\ref{sec:problem}.

The debt framing likewise has a source. Cunningham~\cite{Cunningham1992} introduced technical debt as the deferred cost of an expedient implementation choice, and Sculley et al.~\cite{Sculley2015} extended it to machine-learning systems, where the accumulating liability is configuration, data dependency, and entanglement rather than code. Civic Automation Debt is the same accounting move applied to a public workforce: a locally optimal allocation creates a liability settled later, by someone else, invisible on the ledger that authorized it. Rashidi~\cite{Rashidi2026} carries the concern into the agentic-AI era, arguing that AI is the first technology capable of automating the entry-level cognitive work through which expertise and professional judgment are formed.

\subsection{Fair Division and the Distribution of Scarce Opportunities}
\label{app:rw-fair}

Protected developmental practice is a scarce, valuable, and unequally distributed good, which places the Capability Access Constraint of Section~\ref{sec:access} inside the fair-division literature. Bouveret et al.~\cite{Bouveret2016} survey fair allocation of indivisible goods and the standard solution concepts, proportionality, envy-freeness, and their relaxations, and Caragiannis et al.~\cite{Caragiannis2019} show that maximizing Nash social welfare achieves envy-freeness up to one good together with Pareto optimality.

We use a share constraint against a policy-set target rather than an envy-based criterion, and the reason is institutional rather than technical. Envy-freeness is defined over agents' own valuations of the bundles they receive, whereas a municipal panel must defend a distribution against a statutory equality duty and a published labor-pool composition, which are external standards. A Nash-welfare objective over practitioner valuations would also make the allocation depend on elicited worker preferences that no city currently records and that would be strategically reported if it did. The cost of this choice is that our constraint offers no envy guarantee, which we state plainly in Section~\ref{sec:limitations}.

\subsection{Contestability and Algorithmic Due Process}
\label{app:rw-contest}

A framework that prices civic costs must show its working, or the pricing is unfalsifiable in the way Eubanks~\cite{Eubanks2018} documents for opaque public-sector scoring. Citron~\cite{Citron2008} established the technological-due-process argument that automated administrative decisions require notice, a record, and a route to challenge, and Citron and Pasquale~\cite{CitronPasquale2014} extended it to scoring systems whose inputs the subject cannot inspect. Almada~\cite{Almada2019} shows that a right to human intervention is empty unless the system is designed so a human can actually alter the outcome, and Vaccaro et al.~\cite{Vaccaro2019} document how contestability functions in deployed systems. Alfrink et al.~\cite{Alfrink2023} synthesize this into a design framework of five system features and six development practices, of which the evidentiary record, the intervention request, and the ability to challenge the system's rules rather than one decision are the ones Section~\ref{sec:contest} implements. Moon and Guha~\cite{MoonGuha2026} show that formally fair public-sector prioritization rules can still produce outcomes citizens and caseworkers experience as unjust, which is why the procedure includes a systemic trigger and not only an individual remedy. Article 22 of the General Data Protection Regulation~\cite{GDPR2016} gives the legal floor in EU jurisdictions: a right to obtain human intervention, to express a point of view, and to contest a decision based solely on automated processing.

\subsection{Human--AI Teams, Algorithmic Management, and Ethics}
\label{app:rw-hat}

AI is increasingly a collaborator and a manager rather than an analytical tool. A systematic review of human--AI teams identifies task allocation, communication, interaction, and cognitive augmentation as the major enablers of effective collaboration~\cite{HAT2026}, and algorithmic management already allocates work, issues instructions, monitors performance, and supports managerial decisions~\cite{syed2026fedagent,OECDAM2025}. The OECD notes benefits in consistency, productivity, and decision support alongside concerns about autonomy, privacy, transparency, accountability, discrimination, and safety~\cite{OECD2026Policy}, and Wischnewski et al.~\cite{Wischnewski2023} show these persist even where automation performs well, because miscalibrated operator trust degrades team performance independently of system accuracy. Lipsky~\cite{Lipsky2010} supplies the reason a municipal setting differs: public services are constituted in practice by the discretionary judgments of street-level bureaucrats, so a mechanism that removes discretion does not merely change throughput, it changes what the service is.

Callari et al.~\cite{Callari2024} developed an ethical framework for human--robot collaboration emphasizing autonomy, decision authority, dignity, liability, safety, privacy, and worker resilience. Binding regulation is beginning to formalize parts of this agenda. The EU AI Act~\cite{EUAIAct2024} requires high-risk systems to allow meaningful human oversight by competent, authorized personnel, and the revised ISO~10218 series specifies safety requirements wherever a robot shares space or a task with a human, with Part~1 covering the robot itself~\cite{ISO2025} and Part~2 covering integration, applications, and cells~\cite{ISO2025b}, which is the level at which municipal fleet middleware operates. Such frameworks specify principles or hard physical constraints for designers to respect. They do not supply a computational mechanism translating those principles into real-time decisions about who performs a task.

\subsection{AI, Employment, and Who Bears the Cost}
\label{app:rw-labor}

The employment effect of AI mixes automation, augmentation, task transformation, and organizational change. Autor et al.~\cite{AutorLevyMurnane2003} established the task-based framework on which most estimates rest, and Autor~\cite{Autor2015} argues that displacement is routinely overstated because it ignores the complementarities that create new work. Frey and Osborne~\cite{FreyOsborne2017} estimate occupation-level computerization susceptibility from task composition, and Arntz et al.~\cite{Arntz2016} show that the same evidence analyzed at task rather than occupation level cuts the headline exposure figure by roughly a factor of five. We cite the two together deliberately, because the distance between them is the distance between a policy of pre-emptive displacement management and one of task redesign, and only the second is actionable by an allocation mechanism. Eloundou et al.~\cite{Eloundou2024} extend the task-based approach to large language models, finding far more occupations with at least half their tasks exposed to LLM-enabled automation than to robotics alone, and Acemoglu and Restrepo~\cite{AcemogluRestrepo2020} provide causal rather than exposure-based evidence that robot adoption reduces employment and wages in the most exposed US commuting zones, concentrated among lower-skilled workers.

Exposure is not only occupational. Wood et al.~\cite{Wood2019} document how algorithmic control in the gig economy produces autonomy and precarity simultaneously and unevenly. De Stefano~\cite{DeStefano2016} shows that on-demand and crowdwork arrangements sit largely outside the labor protections municipal employment presumes, and Eubanks~\cite{Eubanks2018} documents how automated public-service decision systems concentrate their errors on the poorest residents. The ILO just-transition guidelines~\cite{ILO2015} set out what an institutional response must contain, namely social dialogue, active labor-market policy, skills anticipation, and social protection, while the OECD~\cite{OECD2026Policy} organizes the same agenda for AI, the World Economic Forum~\cite{WEF2025} finds most employers expecting substantial reskilling needs within five years, and the ILO alongside the United Nations Economic and Social Commission for Western Asia (ESCWA)~\cite{ILOESCWA2026} model Arab-region labor-market transformation to 2035.

Two normative literatures supply our vocabulary for what is preserved. Sen~\cite{Sen1999} and Nussbaum~\cite{Nussbaum2011} define capability as the real freedom to do and to be, which is what a developmental task confers and an automated one withholds. We use the word in that sense, not as a synonym for skill inventory. Ostrom~\cite{Ostrom1990} shows that renewable but depletable shared resources are sustained by monitored, boundary-defined, locally governed appropriation rules rather than by privatization or prohibition, which is the institutional form the budget of Section~\ref{sec:hcpb} takes. What none of these supplies is a decision procedure a city can run on Tuesday morning, for one task, that prices exposure before the allocation is made rather than compensating for it afterwards.

\subsection{Robots, Public Space, and Agentic-City Governance}
\label{app:rw-govern}

Mintrom et al.~\cite{Mintrom2025} find that policies for public-space robots concentrate on safety and privacy while broader social and economic wellbeing receives less attention, motivating anticipatory governance and citizen co-design. Phothong et al.~\cite{Phothong2026} and Nechesov and Ruponen~\cite{NechesovRuponen2024} provide early evidence that structured co-creation and transparent civic-intelligence platforms raise measured trust in AI-enabled municipal decisions, and Aldhi et al.~\cite{Aldhi2025} show that the managerial digital readiness of the municipal workforce itself mediates whether smart-city governance reforms take effect, which is a workforce precondition our framework assumes rather than establishes.

Digital twins are the emerging technical substrate. Temple et al.~\cite{Temple2024} review how city digital twins already assist policy decisions, and D\'iaz-Sarachaga~\cite{DiazSarachaga2025} proposes a governance-assessment framework for them, arguing that institutional and accountability questions have received less attention than technical ones. That is the closest antecedent to the Policy Digital Twin of Section~\ref{sec:policytwin}, and the relationship is worth stating precisely: D\'iaz-Sarachaga~\cite{DiazSarachaga2025} asks whether a deployed urban digital twin is well governed and supplies criteria for assessing it, whereas we ask the converse, whether governance can itself be represented inside the twin as an executable constraint on the decisions the twin makes. Those criteria apply to our Policy Digital Twin as forcefully as to any other component. Ben-Daya et al.~\cite{BenDaya2026} press the same point for critical infrastructure, arguing that a twin used in governance must be able to challenge the assumptions of its own model. Qanazi et al.~\cite{Qanazi2025} review the social dimensions urban twins omit, and the Civic Capability Ledger of Section~\ref{sec:layer2} is one such dimension made computable. On the standards side, the ISO~23247 series supplies a reference architecture for digital twins of production systems and the observable-element modeling it depends on~\cite{Shao2023}, which is the closest published standard for the Urban Task Digital Twin, though it does not cover the policy or workforce elements.

Resilience is related but distinct. Mentges et al.~\cite{Mentges2023} show the term is used inconsistently across the critical-infrastructure literature, while Lian et al.~\cite{Lian2025} show that emergency-response networks recover faster when reassignment after a failure is planned in advance.


\section{Framework Specification Detail}
\label{app:framework}

This appendix holds the specification detail that Section~\ref{sec:framework} states in summary. Nothing here is required to follow the optimization program of Section~\ref{sec:math} or the worked allocation of Section~\ref{sec:worked}; it is the material a municipality implementing the framework would need, and a reviewer checking that each declared quantity is in fact defined.

\subsection{Worker-Group Construction}
\label{app:framework-groups}

Worker groups $\mathcal{G}$ are not the Cartesian product of the protected attributes. That construction yields hundreds of intersectional cells for a mid-size city, nearly all of them empty in a single domain, each carrying a constraint. Instead the Weight Review Panel of Section~\ref{sec:feedback} declares an attribute set $\mathcal{J}$ (occupational class, gender, contract status, district of residence) and a \textit{coarsening} that produces $\mathcal{G}$ subject to a minimum cell size $n^{min}$, with a default of $n^{min}=30$ practitioners per cell per domain. Where the full intersection does not meet $n^{min}$, the panel declares marginal groups, one partition per attribute, which yields $\sum_{j\in\mathcal{J}}|A_j|$ constraints rather than $\prod_{j\in\mathcal{J}}|A_j|$. Attributes may be intersected only where the resulting cells meet $n^{min}$ in the domain concerned.

Cells below $n^{min}$ are reported descriptively with their sample size and are not the subject of a binding constraint. The reason is twofold: a target computed over a handful of practitioners is statistically meaningless, and disclosing per-cell allocation statistics for a handful of identifiable people is itself a data-protection harm, which Section~\ref{sec:dpia} treats. The worked example of Section~\ref{sec:worked} therefore uses a single marginal partition, which is the correct treatment for a domain of 24 practitioners rather than a simplification of a richer one.

\subsection{The Four Remaining Debt Components}
\label{app:framework-dterms}

Section~\ref{sec:framework} states $D^{skill}$ in full as equation~\eqref{eq:dskill}. The remaining four components of equation~\eqref{eq:cad} are each normalized to $[0,1]$ and defined over observables the Civic Capability Ledger and the accountability register already carry.

$D^{fall}_{i,m}$ is the fractional reduction in the fallback capacity $F_{\kappa(i)}$ of equation~\eqref{eq:fallback} that mode $m$ would cause if generalized across the domain. $D^{acct}_{i,m}$ is the share of decision steps lacking a designated accountable human, weighted by the task criticality $crit_i$. $D^{dep}_{i,m}$ is the share of automated inputs drawn from a single vendor with no qualified alternative, following the reasoning by which Sculley et al.~\cite{Sculley2015} treat undeclared dependencies as the dominant hidden liability of deployed machine-learning systems. $D^{trans}_{i,m}$ is the evaluated share of affected workers without an identified reskilling pathway, computed per group before aggregation.

\subsection{Populating the Task Vector}
\label{app:framework-task}

The nine dimensions of equation~\eqref{eq:task} are normalized to $[0,1]$ and populated from sources already used in human--robot collaboration task profiling. Cognitive, physical, and empathy demand ($cog_i$, $phy_i$, $emp_i$) come from standardized task-demand instruments or expert ergonomic rating~\cite{Makrini2019,Merlo2023}. Safety risk $risk_i$ comes from incident history or hazard classification. The legal-authority requirement $auth_i$ and privacy sensitivity $priv_i$ come from the Policy Digital Twin's rule encoding for the sector. Urgency $urg_i$ comes from service-level targets. Human learning value $learn_i$ comes from a structured judgment of developmental value for a junior practitioner, elicited by the protocol of Appendix~\ref{app:elicit}. Criticality $crit_i$ comes from the service's criticality tier in continuity planning.

\subsection{A Worked Policy Encoding}
\label{app:framework-policy}

Article~14 of the EU AI Act~\cite{EUAIAct2024} requires that high-risk systems be effectively overseen by natural persons with the competence, training, and authority to intervene. For a municipal structural-inspection determination the encoding under equation~\eqref{eq:rule} is as follows.

\begin{center}
	\small
	\begin{tabularx}{\textwidth}{@{}lX@{}}
		\toprule
		\textbf{Field} & \textbf{Value}\\
		\midrule
		scope & $\sigma(i)=\textit{structural-inspection} \ \wedge\ m \in \{A,R,A{+}R,H{+}A,H{+}A{+}R\}$\\
		guard & $auth_i \geq 0.8 \ \wedge\ crit_i \geq 0.7$\\
		status & \textit{restrict}, with roster guard $\exists w\in a:\ \mathrm{stage}_k(w,t)=\textit{competent} \wedge \mathrm{license}(w)$\\
		authority & regulation (EU AI Act Art.~14) with national implementing instrument\\
		validity & $[\,2026\text{-}08\text{-}02,\ \infty\,)$\\
		provenance & Panel decision WRP-2026-014 (Workforce Review Panel), encoding reference AIA-14-STRUCT-01 (EU AI Act Art.~14)\\
		\bottomrule
	\end{tabularx}
\end{center}

Under this rule a fully autonomous mode is not prohibited outright. It is restricted to rosters carrying a licensed competent practitioner, which is exactly the difference between a prohibition and an oversight obligation, and is why the two intermediate statuses are kept distinct.

\subsection{The Reskilling Gating Predicate}
\label{app:framework-gating}

Earlier presentations of this framework paired the Just Transition Constraint~\eqref{eq:justtransition} with a reskilling-coverage requirement $\chi_g \geq \chi^{min}$ placed inside the optimization program. That was a category error. The quantity $\chi_g$, the share of affected workers in group $g$ with a funded reskilling or redeployment pathway, is a municipal budget fact and not a function of the allocation $\mathbf{x}$: the optimizer can neither satisfy nor violate it, so stating it as a constraint on $\mathbf{x}$ misrepresents where the obligation sits.

It belongs in the governance layer, and is enforced there as a \textit{gating predicate} evaluated by the Policy Digital Twin. Any staffed mode whose adoption would increase $\Delta_g$ for a group with $\chi_g < \chi^{min}$ receives status \textit{prohibit} until the pathway is funded. Displacement is thereby conditioned on transition support without pretending that the allocator controls the training budget.

\subsection{Policy Conflict Order and Consistency Check}
\label{app:framework-conflict}

Taking the join in equation~\eqref{eq:policy} means the most restrictive applicable rule wins, which is the correct default for a public body. Two rules conflict only when a municipality wishes a lower-ranked rule to \textit{relax} a higher-ranked one, which the twin refuses: a rule may be overridden only by a rule of strictly higher $\mathrm{authority}$, and among equal authority by the later $\mathrm{validity}$ start, with any remaining tie raised to the panel rather than resolved silently. A nightly consistency check reports every scope for which the join is \textit{prohibit} across all modes, since such a scope has no executable allocation and must resolve to Zone~Z4 human-reserved execution or be corrected.

The two intermediate statuses are executable rather than descriptive. \textit{allow-with-oversight} binds the staffed mode to a named reviewer, adds that reviewer's time to $Cost_{i,m,a}$, and raises $\phi_m$ by the reviewer's assessed share of retained judgment. \textit{restrict} narrows the roster set $\mathcal{A}_{i,m}$ to rosters meeting the guard's stated minimum human share or stage requirement. Under the encoding of Appendix~\ref{app:framework-policy}, a fully autonomous mode is therefore restricted to rosters carrying a licensed competent practitioner rather than prohibited outright, which is the difference between a prohibition and an oversight obligation and the reason the two statuses are kept distinct. Algorithm~\ref{alg:scv} implements both.

\subsection{The Four Channels of Developmental Content}
\label{app:framework-channels}

Separating flows from stocks removes double-counting between the objective's terms. It does not mean that developmental content reaches the decision only once. It arrives through four distinct channels: the priced term $w_9\alpha D^{skill}$ in equation~\eqref{eq:scv}, the hard budget~\eqref{eq:hcpb}, the capability dual $\lambda_k$ in equation~\eqref{eq:aug}, and the access dual $\nu_{k,g}$ in the same equation. These are not redundant, because they answer different questions: how much developmental practice is worth, how much must be delivered, what the marginal hour costs, and who receives it. They are nonetheless four channels and not one, and the Lagrangian channel is numerically the largest. Section~\ref{sec:pricediag} makes that comparison explicit through the flip threshold $w_9^\star$, and Section~\ref{sec:pricing} evaluates it for the worked case.

\subsection{The Allocation Market}
\label{app:framework-market}

Ten specialized agents propose and price candidate staffed modes. Five are \textit{capability agents} estimating what a staffed mode can do: the Human Workforce Agent, the AI Capability Broker, the Robot Fleet Agent, the Safety Agent, and the Fiscal and Energy Agent. Five are \textit{civic agents} estimating what it costs the city in stocks it must later replenish: the Workforce Development Agent ($\phi_m$, $\psi_a$, $D^{skill}$), the Equity and Inclusion Agent ($Eq^{srv}$, $D^{trans}$, and the group shares in~\eqref{eq:accessshare}), the Privacy and Rights Agent ($Pr$), the Resilience Agent ($D^{fall}$, $Res^{(n_s)}_s$), and the Policy Guardian Agent, which executes~\eqref{eq:policy} and holds veto authority.

The market protocol is a single-round sealed-bid call. The Policy Guardian Agent publishes the task and the surviving roster set. Each agent returns its terms for every surviving staffed mode within a fixed deadline, with a declared default used where an agent does not answer in time. The auctioneer computes~\eqref{eq:scv} and~\eqref{eq:aug}, applies the admissibility test~\eqref{eq:admis}, and awards. No agent sees another's bid before submitting, so there is no bidding dynamic to converge, and the per-task decision terminates in one round with a bounded number of queries. The terms most vulnerable to strategic reporting are cross-validated against independently observed outcomes, which is a partial response to the principal--agent exposure recorded in Section~\ref{sec:threats}.

\section{Proofs, Normalization Anchors, and Notation}
\label{app:proofs}

This appendix holds the proofs of the four propositions of Section~\ref{sec:architecture}, the anchored normalization of the debt components, the construction of the duals, and the notation table. The proposition statements remain in the main text.

\subsection{Proofs of the Propositions}
\label{app:proofs-props}

\paragraph{Proof of Proposition~\ref{prop:feasibility}.}
	Every term of the left side of~\eqref{eq:hcpb} is bounded above by $\ell_i\bar\zeta_k$, and~\eqref{eq:prog-assign} selects exactly one staffed mode per task, so $\Lambda_k \leq \bar\zeta_k\sum_i \ell_i = \Phi_k\bar\zeta_k$. Substituting~\eqref{eq:hcpb-estimator} and solving for $r_k$ gives~\eqref{eq:rcrit}.

\paragraph{Proof of Proposition~\ref{prop:intake}.}
	A developing practitioner accrues $\ell_i\phi_m\omega_{a,w}=learn_i\,d_i\,\phi_m\,\omega_{a,w}$ qualified-practice hours per task while present for $d_i$ clock hours, hence $\overline{learn}_k\bar\phi_k\bar\omega_k$ hours per clock hour, hence $\Theta_k\overline{learn}_k\bar\phi_k\bar\omega_k$ per period. Dividing $h_k=\overline{learn}_k h^{raw}_k$ from~\eqref{eq:hconv} by that rate cancels $\overline{learn}_k$ and gives~\eqref{eq:tau}. The domain completes $B_k/h_k=\eta_k N_k r_k$ practitioners per period, and by Little's law a queue with that completion rate and residence time $\tau_k$ holds $\eta_k N_k r_k \tau_k$ members.

\paragraph{Proof of Proposition~\ref{prop:price}.}
	Immediate from the definition of the infimum and the continuity of~\eqref{eq:scv} in $w_9$.

\paragraph{Proof of Proposition~\ref{prop:violation}.}
	The resource bracket of~\eqref{eq:admis} rejects any staffed mode whose consumption exceeds the remaining availability, so the running total never exceeds $U_r$ and the bound is exact. For displacement the bracket tests $\Delta_g(t^+)\leq\tau_g$ on the post-execution state, so no accepted task can leave $\Delta_g$ above $\tau_g$, and the tightest bound on the final value is therefore $\tau_g$ itself. Where the estimate $\delta^{disp}_g$ used at decision time is revised after execution by at most $\bar\delta_g$, the realized value can exceed $\tau_g$ by at most that revision, which is one task's worth and does not accumulate because every subsequent task is tested against the updated state.

\subsection{Anchored Normalization of the Debt Components}
\label{app:proofs-anchors}

The five debt components carry no windowed normalization. Each is constructed as a ratio whose denominator is fixed externally, which places each debt term in $[0,1]$ as defined in equation~\eqref{eq:anchors}:
\begin{equation}
	\begin{aligned}
		D^{skill}_{i,m,a} &= 1-\phi_m\psi_a &&\text{against the task's own full developmental content } \ell_i,\\
		D^{fall}_{i,m} &= \min\!\Big(1,\ \Delta F_{\kappa(i)}(m)\,/\,\rho_{\sigma(i)}\Big) &&\text{against the statutory reserve } \rho_s=\kappa_s D^{peak}_s,\\
		D^{acct}_{i,m} &= crit_i \cdot n^{unacct}_{i,m}/n^{steps}_{i,m} &&\text{against the decision-step count of the task},\\
		D^{dep}_{i,m} &= n^{single}_{i,m}/n^{auto}_{i,m} &&\text{against the automated-input count of the mode},\\
		D^{trans}_{i,m} &= \textstyle\sum_{g}\varpi_{g}(i,m)\,\bigl(1-\chi_g\bigr) &&\text{against funded transition coverage } \chi_g,
	\end{aligned}
	\label{eq:anchors}
\end{equation}
where $\Delta F_k(m)$ is the fallback reduction mode $m$ would cause if generalized across the domain and $\varpi_g(i,m)$ is group $g$'s share of the workers whose allocated hours the mode reduces. All five denominators are external to observed performance, so a decade of decline registers as a decade of decline. Note that $D^{trans}$ is anchored to the reskilling coverage $\chi_g$ and not to the displacement ceiling $\tau_g$: the ceiling is a hard constraint on $\Delta_g$, and the priced quantity is the unfunded share of the burden.

\subsection{Duals from a Binary Program}
\label{app:proofs-duals}

Program~\eqref{eq:program} is integral, and integer programs have no exact duals. We take $\lambda,\mu,\nu,\varrho,\upsilon$ from the LP relaxation solved at the root node, which is a valid subgradient of the relaxed value function and the standard construction for Lagrangian decomposition of assignment problems. The consequence is that~\eqref{eq:aug} approximates the relaxed rather than the integral program, and the approximation error is bounded by the integrality gap, which the solver reports at each rebalance and which the protocol of Appendix~\ref{app:protocol} requires to be tabulated at every cadence. In the worked case of Section~\ref{sec:worked} that gap is $7.8\times10^{-6}$ objective units per task, because a single coupling constraint plus the simplex constraint admits at most two active staffed modes at an LP vertex and the domain's task count is large enough that rounding costs almost nothing. We do not claim this holds at every scale, and the reported gap is the diagnostic that says when it stops holding.

\subsection{Notation}
\label{app:proofs-notation}

Table~\ref{tab:notation} summarizes the symbols of Sections~\ref{sec:problem} through~\ref{sec:zones}. Two collisions present in earlier presentations have been removed: the resilience \textit{contribution} term formerly written $Res_{i,m}$ no longer exists, since resilience enters only through $D^{fall}$ and constraint~\eqref{eq:res3r}, leaving $Res^{(n)}_s(t)$ as the unique bearer of that symbol, and safety is written $S_{i,m,a}$ throughout with no separate $Safety$. The single-assignee function of earlier presentations, written inconsistently as $\mathrm{assignee}(i,m)$ and $g(i,m)$, is replaced throughout by the roster group shares $\pi_{a,g}$ of~\eqref{eq:groupshare}.

\begin{table}[!htbp]
	\caption{Notation. Indices: $i$ task, $m$ execution mode, $a$ roster, $w$ practitioner, $k$ capability domain, $s$ critical service, $g$ worker group, $r$ resource, $t$ time.}
	\label{tab:notation}
	\footnotesize
	\begin{tabularx}{\textwidth}{lX}
		\toprule
		\textbf{Symbol} & \textbf{Meaning}\\
		\midrule
		$x_{i,m,a}$ & Binary decision variable: task $i$ executed in mode $m$ with roster $a$, Equation~\eqref{eq:decisionvar}\\
		$\mathcal{M}$; $\mathcal{A}_{i,m}$; $\mathcal{A}(i,t)$ & Mode set; feasible rosters; admissible staffed modes after test~\eqref{eq:admis}\\
		$T_i$; $d_i$ & Task requirement vector and expected duration in hours, Equation~\eqref{eq:task}\\
		$\ell_i$ & Developmental content $learn_i d_i$ in qualified-practice hours, Equation~\eqref{eq:ell}\\
		$c_w(t)$; $\mathrm{stage}_k(w,t)$ & Practitioner $w$'s accrued practice and career stage, Equation~\eqref{eq:stage}\\
		$\omega_{a,w}$; $\psi_a$; $\pi_{a,g}$, $\pi^{dev}_{a,g}$ & Roster work split; developmental eligibility~\eqref{eq:psi}; group shares~\eqref{eq:groupshare}\\
		$\phi_m$ & Share of $\ell_i$ experienced by a human under mode $m$, Equation~\eqref{eq:phi}\\
		$\sigma(i)$, $\kappa(i)$ & Maps from task $i$ to its critical service and capability domain\\
		$L_k(t)$; $F_k(t)$; $E_{k,g}$ & Ledger entry~\eqref{eq:ledger}; fallback capacity~\eqref{eq:fallback}; pipeline by group\\
		$\rho$; $\mathcal{P}(T_i,m,a,t)$ & Policy rule~\eqref{eq:rule}; status join over firing rules~\eqref{eq:policy}\\
		$B_k(t)$; $\Lambda_k(t)$; $\bar B_k(t)$; $\delta_k(t)$ & Budget, accrued practice, pro-rated trajectory, catch-up increment~\eqref{eq:catchup}\\
		$\eta_k, r_k, N_k, N^{dev}_k, h_k, h^{raw}_k, B^{min}_k$ & Replacement factor, attrition, competent and developing headcount, competence hours (weighted and raw), budget floor\\
		$\Phi_k$; $\bar\zeta_k$; $r^{crit}_k$; $\tau_k$; $\hat n_k$ & Developmental content supply, maximum credit, critical attrition~\eqref{eq:rcrit}, formation time~\eqref{eq:tau}, intake requirement~\eqref{eq:intake}\\
		$\Pi_{k,g}$, $\theta_{k,g}$, $\varepsilon_k$, $n^{min}$ & Realized share, target share, tolerance, minimum group cell size\\
		$\Delta_g$, $\tau_g$, $\chi_g$, $\chi^{min}$ & Group displacement, ceiling, funded reskilling coverage and its floor\\
		$Res^{(n)}_s$, $\rho_s$, $\kappa_s$, $D^{peak}_s$, $n_s$, $\xi_{s,k}$ & Surviving capacity~\eqref{eq:res-def}, threshold, demand share, peak demand, contingency order, domain draw share\\
		$Q,S,P,Eq^{srv},Tr,Cost,En,Pr$ & Flow terms of $SCV$, all in $[0,1]$, Equation~\eqref{eq:scv}\\
		$D^{skill},D^{fall},D^{acct},D^{dep},D^{trans}$ & Stock-change components of $CAD$, all in $[0,1]$, Equations~\eqref{eq:cad} and~\eqref{eq:anchors}\\
		$w_1,\ldots,w_9$; $\alpha,\ldots,\epsilon$; $w_9^\star$ & Objective weights, debt weights, and the flip threshold of Proposition~\ref{prop:price}\\
		$\lambda_k,\mu_s,\nu_{k,g},\varrho_g,\upsilon_r$ & Dual prices of~\eqref{eq:prog-hcpb},~\eqref{eq:prog-res},~\eqref{eq:prog-access},~\eqref{eq:prog-just},~\eqref{eq:prog-cap}\\
		$u_{i,m,a,r}$, $U_r$; $S^{min}_i$ & Resource consumption and availability; safety floor $risk_i\,crit_i$\\
		\bottomrule
	\end{tabularx}
\end{table}

\section{Data Protection Impact Assessment Measures}
\label{app:dpia}

Section~\ref{sec:dpia} establishes that the Civic Capability Ledger is high-risk processing under the GDPR~\cite{GDPR2016} and that an Article~35 impact assessment is mandatory before deployment. This appendix sets out the four measures the design assumes that assessment will contain.

\textit{Minimization by aggregation.} The constraints need $\Pi_{k,g}$ and $\Lambda_k$, which are aggregates. Only the career-stage test~\eqref{eq:stage} and the roster construction need per-practitioner data, and both need only the boolean $c_w(t) < h_k$ and the practitioner's availability, not the full practice history. The ledger therefore stores $c_w$ at the granularity certification already requires and exposes only the boolean to the allocator.

\textit{Separation of the protected attribute from the allocator.} Group membership is held by the Equity and Inclusion Agent and returned to the market only as the aggregate share $\pi^{dev}_{a,g}$ for a candidate roster. The optimizer never receives an individual's protected attributes.

\textit{Small-cell suppression.} The $n^{min}$ rule of Section~\ref{sec:problem} is a disclosure control as much as a statistical one. Publishing per-cell practice shares for a cell of three people identifies those three people, which is why cells below $n^{min}$ are reported descriptively with sample sizes and are not the subject of a binding constraint.

\textit{Article 22 and the human decision.} An allocation decided by~\eqref{eq:argmax} without human involvement is a decision based solely on automated processing where it significantly affects the worker, which developmental access and displacement do. The contestability procedure of Section~\ref{sec:contest} supplies the Article 22(3) safeguards, namely the right to obtain human intervention, to express a point of view, and to contest the decision, and the Z3 re-decision remedy is that human intervention in operational form. Zone assignment is therefore not only a safety mechanism but a data-protection one, and a city that places worker-affecting classes in Z1 should expect to defend that under Article 22 rather than only under the AI Act.

\section{Two Parameter Shocks in Full}
\label{app:shocks}

This appendix carries the two parameter shocks of Section~\ref{sec:shocks} through every step of the arithmetic. Every quantity is exact arithmetic on the model of Section~\ref{sec:worked} at the stated parameter values; none is a simulation output.

The sweep of Section~\ref{sec:sensitivity} varies one parameter at a time around the operating point. Two particular displacements are worth carrying through in full, because they are the two regimes Proposition~\ref{prop:feasibility} separates and because a city will meet both. Everything in this subsection is exact arithmetic on the model at the stated parameter values. Nothing here is a simulation output, and the behaviour of an implementation under stochastic arrivals is a separate question that Section~\ref{sec:protocol} specifies and this article does not answer.

\textit{Shock one: a retirement wave inside the feasible region.} Raise attrition from $r_k = 0.12$ to $0.14$, which is inside the critical rate of $0.1418$ at which Proposition~\ref{prop:feasibility} fails for this domain. Equation~\eqref{eq:hcpb-estimator} expands the budget from 4{,}976.64 to $1.2 \times 24 \times 0.14 \times 1440 = 5{,}806.08$ hours, so the required mean credited share rises from 0.2962 to 0.3456. The two active staffed modes are unchanged, so the tie condition of~\eqref{eq:lambdaworked} is unchanged and $\lambda_k$ remains 0.0454: the price of an hour does not move, because the same pair of modes still brackets the requirement. What moves is the mix. The optimum shifts from 28.3\% to 94.1\% robot-assisted human work, since $H{+}R/a_1$ is the only admissible mode whose credited share exceeds the new requirement, and the cost of preservation rises from 4.34\% to 9.56\%. The domain is then operating at $5806.08/5880 = 98.7\%$ of the ceiling Proposition~\ref{prop:feasibility} sets. This is the stabilizing direction the capability feedback is designed to have, and it is worth naming what the arithmetic does \textit{not} establish: that an implementation updating duals at a finite rebalance cadence converges to this point rather than oscillating around it. Proposition~\ref{prop:violation} bounds the drift between rebalances by one task's worth; it does not bound the trajectory, and no result in this article does.

\textit{Shock two: a retirement wave outside it.} Raise attrition to $r_k = 0.25$. The budget becomes $B_k = 10{,}368$ hours against the domain ceiling $\Phi_k\bar\zeta_k = 5{,}880$, so the deficit is 4{,}488 hours, or $4488/1440 = 3.12$ competent-equivalents a year that cannot be formed under \textit{any} allocation of this task stream. The required mean credited share of $10368/16800 = 0.6171$ exceeds both the roster-guard ceiling $\bar\zeta_k = 0.35$ for the admissible modes and the absolute supervision ratio of 0.50 that one developing practitioner per lead implies. Proposition~\ref{prop:feasibility} is violated by arithmetic before any dynamics arise, and the escalation of Section~\ref{sec:feasibility} is what the framework has to offer. Its four options are computable and we state each with its magnitude: permitting two developing practitioners per lead raises $\bar\zeta_k$ to $0.70 \times 2/3 = 0.467$ and leaves a residual deficit of 2{,}528 hours; expanding the supervised stream to $10368/(8 \times 0.35) = 3{,}703$ inspections a year closes it; lowering $\eta_k$ to $5880/(24 \times 0.25 \times 1440) = 0.68$ closes it by deciding not to replace the full cadre, which is a structural contraction and should be recorded as one; and accepting the deficit logs a 3.12-practitioner annual shortfall against a named decision-maker.

The second shock carries the more useful lesson, and it is a lesson about task structure rather than about recruitment. The binding quantity is $\Phi_k\bar\zeta_k$, the practice the task stream can carry, not the number of people a city is willing to hire. A municipality facing this attrition shock could recruit thirty trainees and still fail to avert the shortfall, because the inspection stream supports at most $21{,}000/1{,}500 = 14$ supervised trainee posts at the prevailing automation level. When a capability budget becomes infeasible, the remedy is in the allocation or in the scope of the cadre the city intends to replace. It is not in the intake alone, and a workforce plan that treats it as a hiring problem will spend a recruitment budget and still lose the capability.

\bibliographystyle{Frontiers-Harvard}
\bibliography{references}

\end{document}
